\chapter{Concluzii}

Proiectul de diplomă prezentat ilustrează o metodă de monitorizare a jocurilor de tenis de masă prin vedere artificială, diferită de cele întâlnite în literatura de specialitate. Prin orchestrarea mai multor modele preantrenate descrise în Secțiunea \ref{subsec:static-objects} și Secțiunea \ref{subsec:dynamic-objects} cu proceduri mai simple precum convoluția, eroziunea și dilatarea, se obțin rezultate bune în privința detecției și a performanței. Mai mult, prin introducerea tehnicilor din sfera paralelismului și a concurenței, se valorifică mai eficient echipamentele disponibile.

În comparație cu TTNet, care folosește o rețea neuronală pentru a îndeplini toate cerințele, abordarea noastră asigură o îmbunătățire a performanței de procesare. De asemenea, pentru metricile de detecție a obiectelor, se observă un avans notabil, în ciuda setului de date FMI-TT relativ limitat.

Totuși, proiectul prezintă anumite limitări și posibilități de îmbunătățire. Pentru videoclipurile stereoscopice, ar fi utilă determinarea cu precizie a unghiului dintre camere și transformarea coordonatelor în același sistem de coordonate. Sincronizarea pe baza sunetului poate fi înlocuită fie prin dezvoltarea unei aplicații care să sincronizeze cele două dispozitive de înregistrare, fie prin utilizarea unor metode mai avansate bazate pe histograme audio.

Per ansamblu, considerăm că tehnicile folosite în alcătuirea unui sistem automat de analiză a jocurilor de tenis reprezintă o structură robustă pentru abordarea competițiilor sportive. În același timp, putem spune că scopul proiectului de a identifica și analiza componentele ce alcătuiesc jocurile de tenis este îndeplinit datorită rezultatelor experimentale și vizuale.